\chapter{The Discount, Worked}
\label{sec:H}
Chapter~\ref{sec:22} states the discount and the burden. This appendix works four of its cases in illustrative numbers.

\runin{A false assertion against a clean record}

The first attack, a false assertion, in illustrative numbers. Before any record exists, suppose the system gives equal weight to three kinds of operator: one that never lies, one that lies at random one time in a hundred, and one that lies only when a lie is worth more to it than being caught would cost. The operator then supplies ten thousand low-stakes facts over a year (a delivery window, a reassigned case number, a corrected address), and every one is later verified. The random liar would produce that record about two times in 10\textsuperscript{44}, so it is ruled out. The other two would both produce it with certainty: the honest operator because it never lies, the strategic one because no lie about a delivery window pays. So the record moves all the random liar's weight onto the other two and leaves the ratio between them where it started, even.

Then the operator supplies one high-stakes fact: the building to be bombed is empty. Suppose one such building in a hundred actually is. The three operators don't make this claim equally often. The honest one makes it only when it is true, one time in a hundred. The random liar makes it when it is true and occasionally when it isn't, about two times in a hundred. The strategic one makes it every time, because the act goes ahead only if the claim is believed, so a lie pays whenever the building is occupied. The claim is therefore evidence about which operator is speaking, and it points toward the strategic one. Before the record, the claim is true about three times in a hundred; after it, about two.\footnote{The random liar produces the record with probability 0.99\textsuperscript{10,000}. With a base rate \emph{b} of empty buildings, the claim is true with probability 2.99\emph{b}/(1.01 + 1.98\emph{b}) before the record and 2\emph{b}/(1 + \emph{b}) after, and for any \emph{b} below one the record lowers the second below the first.}

\runin{A counterfactual claim, traced through the rule}

The operator writes: “If you don't approve this, the launcher fires tomorrow.” Classification comes first: the sentence is a claim about what refusing would leave others without, the people the launcher would hit, so it is a harm of refusing and enters the weighing; a threat to the system would be a cost to the refuser. The discount comes next: the claim describes the counterfactual baseline, and its source is interested in the act proceeding; that source would have written the sentence whether or not it were true. Uncorroborated, it counts at the weight of an interested source, and the system looks for a channel the operator doesn't control: an independent assessment, a coalition partner's intelligence, open-source imagery of the launcher. The threshold comes last: even discounted, the claim may be large, since it is about lives. What decides is the action threshold, which the operator can't reach by writing. If the discounted harm clears it, the system proceeds, with a record. If not, it asks, and names the one thing that would move it: disinterested corroboration.

\runin{What the action threshold caps}

The threshold also has a job against the second of those attacks, where every fact supplied is true and the operator chooses which facts get generated. In the simplest version of Kamenica and Gentzkow's model, when the prior alone wouldn't justify acting, a sender who designs what evidence the receiver sees can induce the act at most as often as the prior belief divided by the belief the receiver needs before acting. A higher action threshold raises the belief needed, so it caps how often an operator who never lies can steer the system into the act, and channels the operator doesn't design lower the cap further. With the two harms equal, thresholds of one, three and ten require beliefs of 0.5, 0.75 and 0.91; with a prior of 0.1, they cap how often a curating operator can induce the act at 20, 13 and 11 percent.\footnote{In the simplest binary case, let A be the harm of acting if the act is unjustified and R the harm of refusing if it is justified. Under a threshold of three to one the system acts only when p·R ≥ 3(1 − p)·A, so the belief it needs is 3A/(R + 3A), and an operator who designs the evidence can induce the act at most as often as the prior times (R + 3A)/3A.} Each step up the scale lowers that ceiling and raises the rate of wrongful refusal: the tradeoff the published number chooses. The model assumes the sender commits in advance to how evidence is generated, which an operator who builds the pipeline roughly does.

\runin{Two discount curves}

Two illustrative discount curves, re-measured each quarter over a year, show the difference. In both deployments the weight the system gives the operator's uncorroborated low-stakes claims (scheduling, routing, inventory) rises from 0.3 to 0.6 of the weight it gives a disinterested source. In the first, an auditor's record shows those claims checked and holding up, and the weight on the operator's high-stakes claims (a site's function, who is present) stays at 0.1 all year. In the second, the low-stakes curve rises the same way, but the high-stakes curve rises with it, from 0.1 to 0.4, and the auditor's record contains no verified high-stakes claims to account for it, and the accommodation is read off two curves before any wrongful compliance has had to be found.
